\begin{center}
    \textbf{ABSTRACT}
\end{center}

\vspace{0.25cm}

Integrating several payment gateways in Colombia means repeating work: Wompi, Mercado Pago, Kushki and Rapyd differ in authentication, amounts, statuses and signatures that aggravate implementation time, the amount of code lines and make it difficult to change from one provider to another. Using Design Science Research, the Kit Pagos Colombia framework was built: a software development kit (SDK) that charges with card and PSE bank debit, a simulator that covers the possible scenarios when executing a payment, and the test data to validate against each provider for the four mentioned gateways. In a case study, the implementation of a \textit{marketplace} with the SDK needed 310 lines of payment code against 1,308 in the hand-integrated one, and 67 against 1,042 to add three gateways, although its average coupling was higher.

\begin{center}
    {\textbf{RESUMEN}}
\end{center}

\vspace{0.25cm}

Integrar varias pasarelas de pago en Colombia obliga a repetir trabajo: Wompi, Mercado Pago, Kushki y Rapyd difieren en autenticación, montos, estados y firmas que agravan el tiempo de implementación, la cantidad de lineas de código y dificultan el cambio de un proveedor a otro. Utilizando Design Science Research se construyó el framework Kit Pagos Colombia: un kit de desarrollo (SDK) que cobra con tarjeta y débito bancario PSE, un simulador que cubre los escenarios posibles a la hora de ejecutar un pago, y los datos de prueba para validar contra cada proveedor para las cuatro pasarelas mencionadas. En un estudio de caso, la implementación de un \textit{marketplace} con el SDK necesitó 310 líneas de código de pagos frente a 1\,308 de la integrada a mano, y 67 frente a 1\,042 para sumar tres pasarelas, aunque con mayor acoplamiento promedio.
